\section{Evidence Status and Open Validation}
\label{sec:results}

Table~\ref{tab:evidencemap} (Appendix~\ref{app:artifacts}) summarizes the
full evidence map. Here we state the status of each claim category.

\paragraph{Measured.}
\begin{enumerate}
  \item \emph{Four-arm lifetime benchmark} (Section~\ref{sec:systems},
    Table~\ref{tab:fourarmmeasured}): dense, persistent-low-rank,
    materialized-generated, and fused-no-store arms on Ascend~910B2.
    The shipped fused arm runs at median $0.823$\,ms versus $0.851$\,ms for
    the same-run materialized arm ($3.2\%$ faster) with no core write-back.
    The persistent-low-rank arm remains $3.3\times$ faster than both
    ($0.247$\,ms) at this geometry, with nearly identical persistent
    storage ($13.5$ vs.\ $15.7$\,MiB including the controller); the gap is
    the per-call cost of generating and applying the conditioned core,
    which arms 2 and 4 trade directly---conditioning quality against
    call latency---with no capacity difference.  Arm 2 is therefore the
    honest latency floor for any low-rank arm, and the fused arm's claim
    is narrower: matching the materialized arm while never storing the
    core.
  \item \emph{DDP-synchronized training step} (Appendix~\ref{app:histtables},
    Tables~\ref{tab:ddptrain}~and~\ref{tab:ddptrainopt}): a real
    \texttt{DistributedDataParallel} step on 4$\times$Ascend~910B2 at 8B-scale
    geometry runs $2.87\times$--$3.17\times$ faster than dense with gradient
    all-reduce in the timed region, and $3.37\times$--$4.14\times$ faster with
    the elementwise AdamW update added; the update does not cancel---dense
    pays $8.97$\,ms to update $3dm$ parameters against $1.0$--$1.3$\,ms for
    the low-rank arm's $r(d{+}m)$ parameters.  This is a low-rank
    parameter-count effect on the training side (the core is recomputed in
    backward, not stored); it is distinct from the inference-scoped
    no-store contract.
  \item \emph{Causal multi-seed FFN-layer reconstruction} (Appendix~\ref{app:artifacts},
    Table~\ref{tab:causalquality-app}): four disjoint hash-bound seeds at
    r384, r401, and r512, all rows from one coherent run family. The
    generated diagonal core reaches same-rank parity
    with the trained persistent control at every rank
    (rel.-MSE $0.449$ vs.\ $0.450$ at r384; $0.440$ vs.\ $0.439$ at r401;
    $0.382$ vs.\ $0.383$ at r512). At the near-equal-budget pair
    (persistent r401 at 48.2~MiB vs.\ causal r384 at 48.1~MiB) the
    persistent factor is better by $+1.95\%$ relative MSE---the honest
    equal-budget penalty for the generated arm at this scale.
    Static SVD is markedly worse at every rank ($0.837$ vs.\ $0.450$ at
    r384). A shuffled-context control (Table~\ref{tab:shuffled}) shows that
    context-conditioning contributes only ${\sim}0.05\%$ of MSE
    ($-0.00021$ paired at r384, not detectable at r512); most of the
    benefit comes from the trainable per-token correction itself. A
    rank-512 control (Table~\ref{tab:rank512control}) extends the result:
    same-rank parity persists at r512, while a persistent r512 factor
    beats the causal r384 core by $14.6\%$ relative MSE---capacity, not
    generation, dominates quality.
  \item \emph{Training-lifecycle footprint} (Appendix~\ref{app:artifacts},
    Table~\ref{tab:lifecycle}): recompute-in-backward reduces the
    saved-for-backward activation footprint by $24.5\times$--$71.8\times$
    at exact gradient equality, confirming that ephemeral generation does
    not perturb learning dynamics at the FFN level. A depth-scaling
    measurement (Table~\ref{tab:lifecycle-depth}) confirms the per-layer
    reduction is depth-independent up to a stack of $N{=}16$ layers.
\end{enumerate}

\paragraph{Open.}
The following remain unmeasured or unimplemented:
\begin{enumerate}
  \item \emph{Branch-skipping dispatch}---whole-call routing with a causal
    router and calibrated tail-risk certification
    (Section~\ref{sec:analysis}) is unimplemented. The historical
    \texttt{torch.where} prototype evaluates both branches, so routed
    savings remain unrealized.
  \item \emph{End-to-end inference}---full-model prefill, cached
    single-token decode, and continuous batching with measured off-chip
    bandwidth counters.  (The beyond-L2 epilogue bottleneck itself is now
    resolved at the kernel level: the measured v4 multi-core kernel is
    bit-identical to v1 and up to $3.03\times$ faster at $T{=}4096$,
    Section~\ref{sec:systems}.)
  \item \emph{Fused-kernel training}---the kernel's own training path
    (retained context, in-op generator recomputation, synchronized optimizer
    step); the no-store contract stays inference-only.
\end{enumerate}

The historical single-process conditioning ablation
(Table~\ref{tab:lifetimeablation}, Appendix~\ref{app:histtables}) is
retained as design exploration only; its non-causal controller and
overlapping evaluation invalidate any quality or routing interpretation.
